\documentclass[11pt]{article}

\usepackage[margin=2cm]{geometry}

\usepackage{graphicx}

\usepackage{booktabs}

\usepackage{hyperref}

\usepackage[T1]{fontenc}

\usepackage[utf8]{inputenc}

\title{Auditoría de universos de Santa Ana Xalmimilulco}

\author{Proyecto Atoyac - salida reproducible}

\date{\today}

\begin{document}

\maketitle

\section{Objetivo y jerarquía}

El universo original de Santa Ana es el canónico. MH es una capa ya filtrada de referencia. El foco estricto es una alternativa analítica trazable y el filtrado 2026 es contextual. Ningún conjunto constituye evidencia de descarga o contaminación.

\section{Pipeline único}

La ejecución completa se realiza con \texttt{python scripts\textbackslash run\_santa\_ana\_pipeline.py}: construir alternativa, comparar, documentar y validar.

\section{Tamaños}

\begin{tabular}{ll}
\toprule
conjunto & registros \\
\midrule
universo\_canonico\_original\_santa\_ana & 38 \\
mh\_filtrado & 49 \\
mh\_filtrado\_con\_id & 43 \\
mh\_filtrado\_sin\_id & 6 \\
foco\_estricto\_alternativo & 67 \\
filtrado\_prioritario\_2026 & 42 \\
union\_comparacion & 111 \\
\bottomrule
\end{tabular}

\section{Canónico original frente a MH filtrado}

\begin{tabular}{ll}
\toprule
comparacion\_canonico\_vs\_mh & registros \\
\midrule
canonico\_y\_mh & 7 \\
mh\_sin\_id & 6 \\
solo\_canonico\_original & 31 \\
solo\_mh\_filtrado & 36 \\
\bottomrule
\end{tabular}

\section{Canónico original frente al foco estricto alternativo}

\begin{tabular}{ll}
\toprule
comparacion\_canonico\_vs\_foco & registros \\
\midrule
canonico\_y\_foco\_estricto & 34 \\
solo\_canonico\_original & 4 \\
solo\_foco\_estricto & 33 \\
\bottomrule
\end{tabular}

\section{Reglas de la alternativa analítica}

\begin{tabular}{lll}
\toprule
rule\_id & rule\_type & registros\_activados \\
\midrule
exc\_nombre\_no\_objetivo & exclusion & 35 \\
exc\_actividad\_no\_objetivo & exclusion & 4 \\
exc\_decision\_previa & exclusion & 7 \\
inc\_humedo\_flags & inclusion & 28 \\
inc\_humedo\_texto & inclusion & 34 \\
inc\_humedo\_etapa\_scian & inclusion & 33 \\
inc\_mezclilla\_nombre & inclusion & 42 \\
inc\_mezclilla\_etapa & inclusion & 38 \\
inc\_material\_textil & inclusion & 3 \\
\bottomrule
\end{tabular}

\begin{description}

\item[\texttt{exc\_nombre\_no\_objetivo}] Evita que una palabra generica como lavado o taller se interprete como proceso textil relevante.

\item[\texttt{exc\_actividad\_no\_objetivo}] Retira actividades cuyo giro declarado no corresponde al proceso textil estudiado.

\item[\texttt{exc\_decision\_previa}] Respeta las exclusiones explícitas de la clasificación previa.

\item[\texttt{inc\_humedo\_flags}] Señal estructurada de tratamiento o proceso húmedo textil.

\item[\texttt{inc\_humedo\_texto}] Señal textual explícita de proceso húmedo o acabado.

\item[\texttt{inc\_humedo\_etapa\_scian}] Señal sectorial o de etapa compatible con lavado, tintorería o acabado.

\item[\texttt{inc\_mezclilla\_nombre}] Identifica establecimientos vinculados explícitamente con mezclilla o jeans.

\item[\texttt{inc\_mezclilla\_etapa}] Conserva la señal de mezclilla asignada en la clasificación productiva.

\item[\texttt{inc\_material\_textil}] Incluye producción de material textil, no confección genérica de ropa terminada.

\end{description}

\section{Mapas}

\begin{figure}[h!]\centering

\includegraphics[width=0.95\linewidth]{../outputs/foco_textil_estricto_santa_ana_2026/mapas/mapa_comparacion_directa_4_conjuntos.png}

\caption{Comparación de los cuatro conjuntos.}

\end{figure}

\begin{figure}[h!]\centering

\includegraphics[width=0.95\linewidth]{../outputs/foco_textil_estricto_santa_ana_2026/mapas/mapa_canonico_original_vs_mh_filtrado.png}

\caption{Universo canónico original frente a MH filtrado.}

\end{figure}

\end{document}